% !PRINTED 331
\weberpage{331}

\gesperrt{By the norm} $N(\mathfrak{a})$ of the ideal $\mathfrak{a}$ one understands the
\gesperrt{absolute value} of the determinant $A$ of the substitution $(A)$:
\[
A = \Sigma \pm a_{1,1}\,a_{2,2}\ldots a_{n,n},
\]
and by (3) the norm is independent of the choice of the basis.

1. \gesperrt{We shall call the basis $(\alpha_\nu)$ positive if
the determinant $A$ is positive, and so equal to the norm of $\mathfrak{a}$.
If $(\alpha_\nu)$ is positive, then $(\alpha'_\nu)$ is always
positive when the substitution determinant $L = +1$.}

From every basis one can derive a positive basis by a substitution with
determinant $-1$, for example by an interchange of two of the $\alpha_\nu$, and
in what follows we shall for the most part use only
positive bases.

The basis determined in Bd. II, \S\ 163, (3) \gesperrt{is positive}.

If now $t_1, t_2, \ldots, t_n$ denotes a system of independent variables and
\weq{4}{\lambda = \sum_\nu \alpha_\nu t_\nu}
% Source scan confirms nu as the lower summation index in (4).
is a basis form with positive basis, then there results (Bd. II, \S\ 164)
\weq{5}{N(\lambda) = N(\mathfrak{a})\,T,}
in which
\weq{6}{T = \Sigma \pm t_{1,1}\,t_{2,2}\ldots t_{n,n}}
is a \gesperrt{primitive homogeneous form with integral coefficients}, of degree $n$, in the
variables $t_\nu$. The $t_{r,s}$ are \gesperrt{linear functions} of the $t_\nu$.

If the substitution (2) is applied to (4), there results
\weq{7}{\lambda = \lambda' = \sum \alpha'_\nu t'_\nu,}
if
\weq{8}{(t_1, t_2, \ldots, t_n) = L'(t'_1, t'_2, \ldots, t'_n),}
in which $L'$ is the \gesperrt{transposed substitution} of $L$; and if
$(\alpha'_\nu)$ is likewise a positive basis, then $L'$ has the
determinant $+1$, and formula (5) gives
\weq{9}{N(\lambda) = N(\mathfrak{a})\,T',}
in which $T'$ is a homogeneous function of degree $n$ in the variables $t'$
which arises from $T$ by the substitution (8).

Two forms which arise from one another by a linear substitution with integral
coefficients of determinant $+1$ are called \gesperrt{equivalent}. Sometimes the forms are also
called equivalent when the substitution determinant is $-1$, but then
with the addition ``\gesperrt{improperly}''. It is then proved by (8):
%CONT>>
